% Test fixture for the tone and prior checks of audit_tex.py and for
% reverse_outline.py (first principle; ground rules 10 and 15).
% Planted: machine tone in the abstract and body, a roadmap at the end of the
% introduction, earlier work called a fault, a one-sentence paragraph, a section
% that opens by pointing back, an additive paragraph opener.
% Not flagged: a contrast with earlier work, our own model's "limitation" (the
% limits check covers that), a roadmap-like phrase outside the introduction.
\documentclass{article}
\begin{document}
\begin{abstract}
Here we show a crucial result for antiferromagnetic memory.
\end{abstract}

\section{Introduction}
Dense memory calls for field-free switching of Mn3Sn.
The model of Ref.~\cite{a} fails to explain the threshold current.
Unlike the field-assisted switching of Ref.~\cite{b}, the reversal here needs no field.

Here we show that a tilted kagome plane removes the bias field.
The rest of this paper is organized as follows.

\section{Results}
This tilt sets the threshold current. To model the switching, we solve the equations of motion.
We delve into the reversal dynamics.

Furthermore, the threshold follows the analytical model. It falls with the tilt angle.

\section{Discussion}
In Sec. 2 we describe how the threshold follows the model.

\begin{thebibliography}{2}
\bibitem{a} A. Author, J. Phys. 1, 1 (2020).
\bibitem{b} B. Author, J. Phys. 2, 2 (2021).
\end{thebibliography}
\end{document}
